%% ----------------------------------------------------------------------------
%%  §2.6  Appendice: i numeri complessi
%% ----------------------------------------------------------------------------
\section{Appendice: i numeri complessi}
\label{sec:complessi}

\subsection{Perché ampliare \texorpdfstring{$\R$}{R}}

Consideriamo una generica equazione di secondo grado nell'incognita $z$:
\begin{equation}\label{eq:secondo-grado}
  az^{2}+bz+c=0 \qquad (a\neq0),
\end{equation}
che ha come soluzioni (reali, se $b^{2}-4ac\ge0$)
\begin{equation}\label{eq:formula-risolutiva}
  z_{1,2}=-\frac{b}{2a}\pm\frac{\sqrt{b^{2}-4ac}}{2a}.
\end{equation}
I coefficienti $a$, $b$, $c$ sono legati alle soluzioni $z_{1}$, $z_{2}$ anche dalle
semplici relazioni, che si verificano immediatamente sommando e moltiplicando le
\eqref{eq:formula-risolutiva},
\begin{equation}\label{eq:viete}
  z_{1}+z_{2}=-\frac ba, \qquad z_{1}\cdot z_{2}=\frac ca .
\end{equation}
Proviamo a scrivere tali relazioni nel caso dell'equazione
\begin{equation}\label{eq:z2piu1}
  z^{2}+1=0;
\end{equation}
in tal caso le \eqref{eq:viete} diventano $z_{1}+z_{2}=0$ e $z_{1}\cdot z_{2}=1$.
Queste relazioni, anche se formalmente ben definite, non hanno alcun senso nell'ambito
dei numeri reali, perché non esistono $z_{1}$, $z_{2}$ soluzioni \emph{reali} della
\eqref{eq:z2piu1}: il quadrato di un numero reale non è mai negativo. Questo esempio
mostra come talvolta sia utile pensare che qualsiasi equazione di secondo grado ammetta
soluzione. Dato che ciò non vale nell'ambito dei numeri reali, si amplia $\R$
introducendo il campo $\C$ dei \emph{numeri complessi}.

\begin{notastorica}[Cardano, Bombelli e il «più di meno»]
Curiosamente, i numeri complessi non nacquero dalle equazioni di secondo grado (di fronte
a $z^{2}+1=0$ ci si limitava a dire che non ci sono soluzioni) ma da quelle di
\emph{terzo} grado. Nell'\emph{Ars Magna} (1545) Gerolamo Cardano pubblicò la formula
risolutiva dell'equazione $x^{3}=px+q$, scoperta da Scipione del Ferro e da Tartaglia.
Applicata a $x^{3}=15x+4$, che ha la soluzione reale $x=4$, la formula produce
\[
  x=\sqrt[3]{2+\sqrt{-121}}+\sqrt[3]{2-\sqrt{-121}},
\]
un'espressione con radici di numeri negativi. Fu il bolognese Rafael Bombelli, nella sua
\emph{L'Algebra} (1572), ad avere il coraggio di calcolare con questi oggetti
«sofistici», introducendo regole come «più di meno via più di meno fa meno» (cioè
$i\cdot i=-1$). Bombelli osservò che $(2+\sqrt{-1})^{3}=2+\sqrt{-121}$ e
$(2-\sqrt{-1})^{3}=2-\sqrt{-121}$, e quindi
$x=(2+\sqrt{-1})+(2-\sqrt{-1})=4$: i numeri «impossibili» servivano a trovare una
soluzione perfettamente reale. Il lettore potrà verificare il calcolo di Bombelli dopo
aver letto questo paragrafo.
\end{notastorica}

\subsection{Definizione e forma algebrica}

Il testo di riferimento \cite{marcellini-sbordone} introduce i numeri complessi
definendo «$i=\sqrt{-1}$» e calcolando «con le stesse regole dei reali». È un approccio
efficace, ma lascia aperta una domanda: \emph{esiste} davvero un sistema di numeri che
contiene $\R$ e un elemento con quadrato $-1$, nel quale valgano tutte le regole del
calcolo? La risposta più semplice, dovuta a William Rowan Hamilton (1837), consiste nel
costruirlo esplicitamente come insieme di coppie di numeri reali.

\needspace{8\baselineskip}
\begin{definizione}[Numeri complessi]\index{numeri complessi}
Si chiama insieme dei \emph{numeri complessi}, e si indica con $\C$, l'insieme
$\R^{2}=\R\times\R$ delle coppie ordinate di numeri reali, munito delle operazioni
\begin{align}
  (x,y)+(x',y')&=(x+x',\ y+y'), \label{eq:somma-complessi}\\
  (x,y)\cdot(x',y')&=(xx'-yy',\ xy'+x'y). \label{eq:prodotto-complessi}
\end{align}
\end{definizione}

La somma è quella «componente per componente»; il prodotto, a prima vista misterioso, è
scelto apposta perché funzioni la costruzione che segue.

\begin{teorema}[$\C$ è un campo]\label{teo:C-campo}\index{campo!dei numeri complessi}
Con le operazioni \eqref{eq:somma-complessi}, \eqref{eq:prodotto-complessi} l'insieme
$\C$ soddisfa gli assiomi (A1)--(A6) del \cref{sec:assiomi}, con elementi neutri
$(0,0)$ e $(1,0)$, opposto $-(x,y)=(-x,-y)$ e, per $(x,y)\neq(0,0)$, inverso
\begin{equation}\label{eq:inverso-complesso}
  (x,y)^{-1}=\Bigl(\frac{x}{x^{2}+y^{2}},\ \frac{-y}{x^{2}+y^{2}}\Bigr).
\end{equation}
\end{teorema}

\begin{schema}
\Ipotesi{$\C=\R^{2}$ con le operazioni \eqref{eq:somma-complessi},
\eqref{eq:prodotto-complessi}; \ gli assiomi di campo valgono in $\R$.}
\Tesi{$\C$ soddisfa (A1)--(A6), con zero $(0,0)$, unità $(1,0)$, opposto $(-x,-y)$ e,
per $(x,y)\neq(0,0)$, inverso dato dalla \eqref{eq:inverso-complesso}.}
\end{schema}

\begin{strategia}
Si tratta di verificare uno per uno gli assiomi, usando le proprietà dei numeri reali
sulle componenti. Le proprietà della somma sono immediate (si lavora componente per
componente). Per il prodotto, la commutatività si vede dalla simmetria della formula;
l'associatività e la distributiva richiedono un calcolo diretto, che svolgiamo; per
l'inverso basta verificare che la coppia proposta funziona, dopo aver osservato che il
denominatore $x^{2}+y^{2}$ non si annulla.
\end{strategia}

\begin{dimostrazione}
Le proprietà associativa e commutativa della somma, l'esistenza dello zero $(0,0)$ e
dell'opposto $(-x,-y)$ seguono immediatamente dalle corrispondenti proprietà dei numeri
reali, applicate alle due componenti. Il prodotto è commutativo, perché scambiando
$(x,y)$ con $(x',y')$ la \eqref{eq:prodotto-complessi} non cambia. È associativo: posto
$z=(a,b)$, $w=(c,d)$, $u=(e,f)$,
\begin{align*}
  (zw)u&=(ac-bd,\ ad+bc)\cdot(e,f)\\
       &=(ace-bde-adf-bcf,\ acf-bdf+ade+bce),\\
  z(wu)&=(a,b)\cdot(ce-df,\ cf+de)\\
       &=(ace-adf-bcf-bde,\ acf+ade+bce-bdf),
\end{align*}
e le due coppie coincidono. È distributivo:
\begin{align*}
  (a,b)\cdot\bigl[(c,d)+(e,f)\bigr]&=\bigl(a(c+e)-b(d+f),\ a(d+f)+b(c+e)\bigr)\\
  &=(ac-bd,\ ad+bc)+(ae-bf,\ af+be).
\end{align*}
L'elemento $(1,0)$ è neutro: $(x,y)(1,0)=(x\cdot1-y\cdot0,\ x\cdot0+1\cdot y)=(x,y)$; ed
è diverso da $(0,0)$. Infine, se $(x,y)\neq(0,0)$, allora $x^{2}+y^{2}>0$ (una somma di
quadrati di numeri reali è nulla solo se entrambi sono nulli, perché ogni quadrato è
$\ge0$), e
\[
  (x,y)\cdot\Bigl(\frac{x}{x^{2}+y^{2}},\frac{-y}{x^{2}+y^{2}}\Bigr)
  =\Bigl(\frac{x^{2}+y^{2}}{x^{2}+y^{2}},\ \frac{-xy+xy}{x^{2}+y^{2}}\Bigr)=(1,0). \qedhere
\]
\end{dimostrazione}

Poiché valgono gli assiomi (A1)--(A6), in $\C$ valgono anche \emph{tutte} le loro
conseguenze dimostrate nel \cref{sec:conseguenze}: leggi di semplificazione, legge di
annullamento del prodotto, regola dei segni, unicità di opposto e inverso. Ecco il
vantaggio dell'approccio assiomatico: non serve ridimostrare nulla.

\paragraph{I reali dentro i complessi.} Le coppie della forma $(x,0)$ si sommano e si
moltiplicano esattamente come i numeri reali:
\[
  (x,0)+(x',0)=(x+x',0), \qquad (x,0)\cdot(x',0)=(xx',0).
\]
Possiamo quindi \emph{identificare} il numero reale $x$ con la coppia $(x,0)$, e
considerare $\R$ come un sottoinsieme di $\C$. Poniamo poi
\[
  i=(0,1),
\]
e risulta
\begin{equation}\label{eq:unita-immaginaria}
  i^{2}=(0,1)\cdot(0,1)=(0\cdot0-1\cdot1,\ 0\cdot1+1\cdot0)=(-1,0)=-1.
\end{equation}
Il numero $i$ si chiama \emph{unità immaginaria}\index{unità immaginaria}. Infine, per
ogni coppia,
\[
  (x,y)=(x,0)+(0,y)=(x,0)+(y,0)\cdot(0,1)=x+iy .
\]
Abbiamo così ottenuto la \emph{forma algebrica} dei numeri complessi:
\begin{equation}\label{eq:forma-algebrica}
  \C=\{z=x+iy : x,y\in\R\}.
\end{equation}
Si dice che $x=\operatorname{Re}z$ è la \emph{parte reale} e $y=\operatorname{Im}z$ è
il \emph{coefficiente della parte immaginaria} di $z$. Due numeri complessi sono uguali
se e solo se hanno uguali parte reale e coefficiente dell'immaginario (è la
\eqref{eq:coppie}). Le operazioni sui numeri complessi si eseguono con le stesse regole
dei numeri reali, tenendo anche conto della \eqref{eq:unita-immaginaria}; quindi:
\begin{align}
  (x+iy)+(x'+iy')&=(x+x')+i(y+y'); \label{eq:somma-algebrica}\\
  (x+iy)\cdot(x'+iy')&=xx'+ixy'+ix'y+i^{2}yy' \notag\\
                      &=(xx'-yy')+i(xy'+x'y). \label{eq:prodotto-algebrico}
\end{align}
La \eqref{eq:prodotto-algebrico} spiega, a posteriori, la scelta del prodotto
\eqref{eq:prodotto-complessi}: è l'unica compatibile con la proprietà distributiva e con
$i^{2}=-1$.

\begin{notastorica}[Hamilton: dalle coppie ai quaternioni]
William Rowan Hamilton (1805--1865) pubblicò la sua teoria delle «coppie algebriche» nel
1837, negli atti dell'Accademia irlandese: voleva fondare i numeri complessi sulla sola
algebra, senza bisogno del piano di Wessel e di Argand di cui diremo più avanti. Il
successo lo spinse a cercare per lo \emph{spazio} ciò che le coppie fanno per il piano:
una moltiplicazione di terne $(x,y,z)$ che rappresentasse le rotazioni tridimensionali,
come quella dei complessi rappresenta le rotazioni piane. Non poteva riuscirci (oggi
sappiamo che in dimensione $3$ una moltiplicazione con le proprietà richieste non esiste),
ma il 16 ottobre 1843, camminando lungo il Royal Canal di Dublino, capì che bastava
passare alle \emph{quaterne}, con tre unità immaginarie legate dalle relazioni
$i^{2}=j^{2}=k^{2}=ijk=-1$, che incise sulla pietra del Brougham Bridge. Il prezzo era la
rinuncia alla commutatività: $ij=k$, ma $ji=-k$. I \emph{quaternioni}, a lungo considerati
una curiosità, sono oggi di uso quotidiano in ingegneria aerospaziale: i sistemi di
navigazione inerziale e di controllo d'assetto di velivoli e veicoli spaziali descrivono
l'orientamento con un quaternione di modulo $1$, che non soffre delle singolarità degli
angoli di Eulero.
\end{notastorica}

\begin{proposizione}[$\C$ non è ordinabile]\label{prop:C-non-ordinabile}
Non esiste in $\C$ alcuna relazione $\le$ che soddisfi gli assiomi dell'ordinamento
(O1)--(O4).
\end{proposizione}

\begin{schema}
\Ipotesi{$\C$ è un campo nel quale $i^{2}=-1$.}
\Tesi{Non esiste in $\C$ una relazione $\le$ che verifichi (O1)--(O4).}
\end{schema}

\begin{strategia}
Supponiamo per assurdo che un tale ordinamento esista. Allora valgono in $\C$ tutte le
conseguenze degli assiomi del \cref{sec:conseguenze}, e in particolare un fatto che si
deduce dai soli assiomi (lo verifichiamo nella dimostrazione): \emph{ogni quadrato è $\ge0$}.
Ma in $\C$ anche $-1$ è un quadrato, $-1=i^{2}$; e $1=1^{2}$ è un quadrato. Avremmo
dunque $1\ge0$ e $-1\ge0$, da cui $1=0$.
\end{strategia}

\begin{dimostrazione}
Se esistesse un tale ordinamento, per ogni $z\in\C$ si avrebbe $z^{2}\ge0$: se $z\ge0$
per l'assioma \ref{ass:ordine-prodotto}; se $z\le0$, allora $-z\ge0$
(\cref{prop:segno-opposto}) e $z^{2}=(-z)^{2}\ge0$ (\cref{prop:segni}). In particolare
$1=1^{2}\ge0$ e $-1=i^{2}\ge0$; dalla seconda, per la \cref{prop:segno-opposto},
$1\le0$, e per l'antisimmetria \ref{ass:antisimmetrica} $1=0$, contro l'assioma
\ref{ass:neutri}.
\end{dimostrazione}

\begin{attenzione}[Nei complessi non si scrivono disuguaglianze]
Una conseguenza pratica: scrivere $z<w$ con $z,w$ complessi non reali non ha senso. Le
disuguaglianze si scrivono solo tra numeri reali, ad esempio tra i \emph{moduli}
$\abs z$, $\abs w$, che definiamo tra poco.
\end{attenzione}

\subsection{Coniugato e modulo}

\begin{definizione}[Coniugato e modulo]\index{complesso coniugato}\index{modulo!di un numero complesso}
Il numero complesso $\bar z=x-iy$ si chiama \emph{complesso coniugato} del numero
$z=x+iy$. Il numero reale non negativo
\begin{equation}\label{eq:modulo-complesso}
  \abs z=\sqrt{x^{2}+y^{2}}
\end{equation}
si chiama \emph{modulo} di $z$. Se $z=x$ è reale, $\abs z=\sqrt{x^{2}}$ coincide con il
valore assoluto di $x$: la notazione è coerente.
\end{definizione}

\begin{proposizione}[Proprietà del coniugato e del modulo]\label{prop:coniugato-modulo}
Per ogni $z,w\in\C$:
\begin{enumerate}
  \item $\overline{z+w}=\bar z+\bar w$, \quad $\overline{zw}=\bar z\,\bar w$, \quad
  $\bar{\bar z}=z$;
  \item $z\,\bar z=\abs z^{2}$; \quad $\operatorname{Re}z=\dfrac{z+\bar z}{2}$, \quad
  $\operatorname{Im}z=\dfrac{z-\bar z}{2i}$;
  \item $\abs{zw}=\abs z\,\abs w$; \quad $\abs{\bar z}=\abs z$; \quad
  $\abs{\operatorname{Re}z}\le\abs z$, \quad $\abs{\operatorname{Im}z}\le\abs z$;
  \item $\abs{z+w}\le\abs z+\abs w$ \quad (disuguaglianza triangolare).
\end{enumerate}
\end{proposizione}

\begin{schema}
\Ipotesi{$z=x+iy$, \ $w=u+iv$, \ con $x,y,u,v\in\R$.}
\Tesi{Le identità (i)--(ii), le relazioni (iii) sui moduli e la disuguaglianza
triangolare $\abs{z+w}\le\abs z+\abs w$.}
\end{schema}

\begin{strategia}
(i) e (ii) sono calcoli diretti con la forma algebrica. Il punto (ii),
$z\bar z=\abs z^{2}$, è lo strumento principale per il resto: permette di trasformare
affermazioni sui moduli (che contengono radici) in affermazioni sui \emph{quadrati} dei
moduli, che sono prodotti e si calcolano con l'algebra. Così per (iii) calcoleremo
$\abs{zw}^{2}$ e per la disuguaglianza triangolare $\abs{z+w}^{2}$; alla fine si
torna ai moduli estraendo la radice, che è crescente (\cref{cor:radice-2}).
\end{strategia}

\begin{dimostrazione}
(i) Se $z=x+iy$, $w=u+iv$, si ha $zw=(xu-yv)+i(xv+yu)$ e
$\bar z\bar w=(x-iy)(u-iv)=(xu-yv)-i(xv+yu)=\overline{zw}$; le altre due sono immediate.
(ii) $z\bar z=(x+iy)(x-iy)=x^{2}-i^{2}y^{2}=x^{2}+y^{2}=\abs z^{2}$; inoltre
$z+\bar z=2x$ e $z-\bar z=2iy$.
(iii) Per (i) e (ii), $\abs{zw}^{2}=zw\,\overline{zw}=z\bar z\,w\bar w=\abs z^{2}\abs w^{2}$,
e poiché i moduli sono non negativi, $\abs{zw}=\abs z\abs w$. La seconda è ovvia; per la
terza, $\abs{\operatorname{Re}z}=\sqrt{x^{2}}\le\sqrt{x^{2}+y^{2}}$ per la monotonia della
radice, e analogamente per la quarta.
(iv) Usando (i), (ii) e il fatto che $\bar zw=\overline{z\bar w}$,
\[
  \abs{z+w}^{2}=(z+w)(\bar z+\bar w)=\abs z^{2}+\abs w^{2}+z\bar w+\overline{z\bar w}
  =\abs z^{2}+\abs w^{2}+2\operatorname{Re}(z\bar w).
\]
Per (iii), $\operatorname{Re}(z\bar w)\le\abs{z\bar w}=\abs z\abs w$, quindi
$\abs{z+w}^{2}\le\abs z^{2}+\abs w^{2}+2\abs z\abs w=(\abs z+\abs w)^{2}$; estraendo la
radice si ottiene la tesi.
\end{dimostrazione}

Il coniugato è particolarmente utile nel calcolo del quoziente di due numeri complessi:
moltiplicando numeratore e denominatore per il coniugato del denominatore, il
denominatore diventa reale. Se $z=x+iy\neq0$,
\begin{equation}\label{eq:quoziente-algebrico}
\begin{split}
  \frac{x'+iy'}{x+iy}&=\frac{(x'+iy')(x-iy)}{(x+iy)(x-iy)}=\frac{(x'+iy')(x-iy)}{x^{2}+y^{2}}\\
  &=\frac{xx'+yy'}{x^{2}+y^{2}}-i\,\frac{x'y-xy'}{x^{2}+y^{2}}.
\end{split}
\end{equation}
Come si vede dal conto precedente, la divisione tra due numeri complessi è possibile
purché il denominatore $z=x+iy$ sia diverso da zero, cioè purché non risulti
contemporaneamente $x=0$ e $y=0$ (è la \eqref{eq:inverso-complesso}).

\begin{esempio}
Calcoliamo un quoziente moltiplicando per il coniugato del denominatore:
\[
  \frac{3+2i}{1-i}=\frac{(3+2i)(1+i)}{(1-i)(1+i)}=\frac{3+3i+2i+2i^{2}}{2}=\frac{1+5i}{2}
  =\frac12+\frac52\,i .
\]
\end{esempio}

\subsection{Il piano di Gauss e la forma trigonometrica}

È utile la \emph{rappresentazione cartesiana} dei numeri complessi, che si ottiene
facendo corrispondere a ogni numero complesso $z=x+iy$ il punto $P$ di coordinate
$(x,y)$, come nella \cref{fig:piano-gauss}: l'asse delle ascisse rappresenta i numeri
reali (\emph{asse reale}), quello delle ordinate i multipli reali di $i$ (\emph{asse
immaginario}). Il piano così interpretato si chiama \emph{piano di Gauss} (o di
Argand--Gauss). Il modulo $\abs z$ è la distanza di $P$ dall'origine (per il teorema di
Pitagora), il coniugato $\bar z$ è il simmetrico di $z$ rispetto all'asse reale, e la
somma di due numeri complessi corrisponde alla regola del parallelogramma per la somma
di vettori.

\begin{figure}[htb]
\centering
\begin{tikzpicture}[scale=1.05, every node/.style={font=\footnotesize}]
  \draw[-Latex, grigio] (-0.5,0) -- (4.3,0) node[below] {$\operatorname{Re}$};
  \draw[-Latex, grigio] (0,-2.3) -- (0,2.6) node[left] {$\operatorname{Im}$};
  \coordinate (z) at (3,1.8);
  \draw[blu, thick, -Latex] (0,0) -- (z) node[midway, above left] {$\rho=\abs z$};
  \fill[blu] (z) circle (1.6pt) node[above right] {$z=x+iy$};
  \draw[tratteggio] (z) -- (3,0) node[below] {$x$};
  \draw[tratteggio] (z) -- (0,1.8) node[left] {$y$};
  \draw[rame, thick] (0.8,0) arc (0:31:0.8) node[midway, right] {$\vartheta$};
  \coordinate (zb) at (3,-1.8);
  \draw[verde, thick, -Latex] (0,0) -- (zb);
  \fill[verde] (zb) circle (1.6pt) node[below right] {$\bar z=x-iy$};
  \draw[tratteggio] (3,0) -- (zb);
  \node[below left] at (0,0) {$0$};
\end{tikzpicture}\hfill
\begin{tikzpicture}[scale=1.05, every node/.style={font=\footnotesize}]
  \draw[-Latex, grigio] (-0.5,0) -- (4.3,0) node[below] {$\operatorname{Re}$};
  \draw[-Latex, grigio] (0,-0.5) -- (0,3.3) node[left] {$\operatorname{Im}$};
  \coordinate (z) at (2.6,0.8); \coordinate (w) at (0.8,2.0); \coordinate (s) at (3.4,2.8);
  \fill[bluchiaro] (0,0) -- (z) -- (s) -- (w) -- cycle;
  \draw[blu, thick, -Latex] (0,0) -- (z) node[below right] {$z$};
  \draw[rame, thick, -Latex] (0,0) -- (w) node[above left] {$w$};
  \draw[verde, very thick, -Latex] (0,0) -- (s) node[above right] {$z+w$};
  \draw[tratteggio] (z) -- (s) -- (w);
\end{tikzpicture}
\caption{A sinistra: il piano di Gauss; il numero complesso $z$ è il punto di coordinate
$(x,y)$, individuato anche dal modulo $\rho$ e dall'argomento $\vartheta$; il coniugato
$\bar z$ è il simmetrico rispetto all'asse reale. A destra: la somma di numeri complessi
segue la regola del parallelogramma; la disuguaglianza triangolare dice che la diagonale
non supera la somma dei lati.}
\label{fig:piano-gauss}
\end{figure}

Un punto $P\neq0$ del piano può essere individuato anche dalla distanza $\rho$
dall'origine e dall'angolo $\vartheta$ che il segmento $OP$ forma con il semiasse reale
positivo. Dal teorema di Pitagora si deduce che
\begin{equation}\label{eq:rho}
  \rho=\sqrt{x^{2}+y^{2}}=\abs z;
\end{equation}
l'angolo $\vartheta$ è legato alla parte reale $x$ e al coefficiente dell'immaginario
$y$ dalle formule
\begin{equation}\label{eq:theta}
  \cos\vartheta=\frac x\rho, \qquad \sen\vartheta=\frac y\rho .
\end{equation}
Un tale $\vartheta$ esiste: poiché $(x/\rho)^{2}+(y/\rho)^{2}=1$, il punto
$(x/\rho,y/\rho)$ sta sulla circonferenza di raggio $1$, e per la definizione di seno e
coseno (\cref{sec:trigonometriche}) è il punto corrispondente a un opportuno angolo
$\vartheta$.

\begin{definizione}[Argomento e forma trigonometrica]\index{argomento di un numero complesso}\index{forma trigonometrica}
Se $z\neq0$, il numero $\vartheta$ della \eqref{eq:theta}, misurato in radianti, si
chiama \emph{argomento} di $z$; esso è definito a meno di multipli interi di $2\pi$.
Tenendo conto delle relazioni tra $\rho$, $\vartheta$ e $x$, $y$, il numero complesso
$z=x+iy$ si scrive in \emph{forma trigonometrica}
\begin{equation}\label{eq:forma-trigonometrica}
  z=\rho(\cos\vartheta+i\sen\vartheta).
\end{equation}
\end{definizione}

\begin{teorema}[Prodotto e quoziente in forma trigonometrica]\label{teo:prodotto-trigonometrico}
Siano
\begin{equation}\label{eq:due-complessi}
  z=\rho(\cos\vartheta+i\sen\vartheta), \qquad z'=\rho'(\cos\vartheta'+i\sen\vartheta').
\end{equation}
Allora
\begin{equation}\label{eq:prodotto-trigonometrico}
  z\cdot z'=\rho\rho'\bigl[\cos(\vartheta+\vartheta')+i\sen(\vartheta+\vartheta')\bigr],
\end{equation}
e, se $z'\neq0$,
\begin{equation}\label{eq:quoziente-trigonometrico}
  \frac{z}{z'}=\frac{\rho}{\rho'}\bigl[\cos(\vartheta-\vartheta')+i\sen(\vartheta-\vartheta')\bigr].
\end{equation}
Cioè: \emph{il prodotto di due numeri complessi ha per modulo il prodotto dei moduli e
per argomento la somma degli argomenti}; il quoziente ha per modulo il quoziente dei
moduli e per argomento la differenza degli argomenti.
\end{teorema}

\begin{schema}
\Ipotesi{$z$, $z'$ in forma trigonometrica \eqref{eq:due-complessi}, con
$\rho,\rho'\ge0$; \ per il quoziente, $z'\neq0$ (cioè $\rho'>0$).}
\Tesi{Valgono la \eqref{eq:prodotto-trigonometrico} e la
\eqref{eq:quoziente-trigonometrico}.}
\end{schema}

\begin{strategia}
Si moltiplica in forma algebrica e si riconoscono, nella parte reale e nel coefficiente
dell'immaginario, i secondi membri delle formule di addizione del coseno e del seno
(\eqref{eq:addizione-cos}, \eqref{eq:addizione-sen}). Per il quoziente conviene non
rifare il conto: si calcola l'inverso di $\cos\vartheta'+i\sen\vartheta'$, che è il suo
coniugato $\cos(-\vartheta')+i\sen(-\vartheta')$, e si applica la formula del prodotto.
\end{strategia}

\begin{dimostrazione}
Tenendo conto delle formule di addizione \eqref{eq:addizione-sen},
\eqref{eq:addizione-cos} si ottiene
\begin{equation}\label{eq:prodotto-passo}
\begin{split}
  z\cdot z'&=\rho\rho'(\cos\vartheta+i\sen\vartheta)(\cos\vartheta'+i\sen\vartheta')\\
  &=\rho\rho'\bigl[(\cos\vartheta\cos\vartheta'-\sen\vartheta\sen\vartheta')\\
  &\qquad\quad+i(\sen\vartheta\cos\vartheta'+\sen\vartheta'\cos\vartheta)\bigr]\\
  &=\rho\rho'\bigl[\cos(\vartheta+\vartheta')+i\sen(\vartheta+\vartheta')\bigr].
\end{split}
\end{equation}
Per il quoziente, osserviamo che
$(\cos\vartheta'+i\sen\vartheta')(\cos\vartheta'-i\sen\vartheta')=\cos^{2}\vartheta'+\sen^{2}\vartheta'=1$
per la relazione fondamentale \eqref{eq:relazione-fondamentale}; quindi, essendo il
coseno pari e il seno dispari,
\[
  \frac{1}{z'}=\frac{1}{\rho'}(\cos\vartheta'-i\sen\vartheta')
  =\frac{1}{\rho'}\bigl[\cos(-\vartheta')+i\sen(-\vartheta')\bigr],
\]
e la \eqref{eq:quoziente-trigonometrico} segue dalla \eqref{eq:prodotto-trigonometrico}
applicata a $z\cdot\frac{1}{z'}$.
\end{dimostrazione}

\begin{figure}[htb]
\centering
\begin{tikzpicture}[scale=1.25, every node/.style={font=\footnotesize}]
  \draw[-Latex, grigio] (-0.4,0) -- (3.6,0) node[below] {$\operatorname{Re}$};
  \draw[-Latex, grigio] (0,-0.3) -- (0,3.2) node[left] {$\operatorname{Im}$};
  \draw[grigio!50] (1,0) arc (0:100:1);
  % z: modulo 1.6, arg 20°; w: modulo 1.25, arg 45°; zw: modulo 2, arg 65°
  \draw[blu, thick, -Latex] (0,0) -- (20:1.6) node[right] {$z$};
  \draw[rame, thick, -Latex] (0,0) -- (45:1.25) node[above right] {$z'$};
  \draw[verde, very thick, -Latex] (0,0) -- (65:2.0) node[above] {$zz'$};
  \draw[blu] (0.55,0) arc (0:20:0.55); \node[blu] at (10:0.75) {$\vartheta$};
  \draw[verde] (20:0.95) arc (20:65:0.95); \node[verde] at (32:1.2) {$\vartheta'$};
  \node[grigio, align=left, anchor=west] at (1.9,2.6) {$\abs{zz'}=\abs z\,\abs{z'}$\\ $\arg(zz')=\vartheta+\vartheta'$};
\end{tikzpicture}
\caption{Il significato geometrico della \eqref{eq:prodotto-trigonometrico}: moltiplicare
per $z'$ significa \emph{ruotare} di un angolo $\vartheta'$ e \emph{dilatare} di un
fattore $\abs{z'}$. In particolare, moltiplicare per un numero di modulo $1$ è una
rotazione pura; moltiplicare per $i$ ruota di un angolo retto.}
\label{fig:prodotto-rotazione}
\end{figure}

\begin{notastorica}[Wessel, Argand, Gauss: i numeri diventano punti]
Per più di due secoli dopo Bombelli i numeri complessi furono usati con profitto, ma senza
che nessuno sapesse dire che cosa fossero: ancora Eulero, che li maneggiava con grande
disinvoltura, li chiamava «impossibili». La svolta venne dalla geometria. Il norvegese
Caspar Wessel (1745--1818), cartografo al servizio dell'Accademia danese delle scienze,
cercava un modo di calcolare con i segmenti orientati: in una memoria presentata
all'Accademia nel 1797, e pubblicata nel 1799, li sommava con la regola del
parallelogramma e li moltiplicava moltiplicando le lunghezze e sommando gli angoli,
proprio come nel \cref{teo:prodotto-trigonometrico}. Scritta in danese, la memoria restò
ignorata fino alla traduzione francese del 1897. Nel 1806 Jean-Robert Argand, un
dilettante che viveva a Parigi, fece circolare in poche copie, senza firmarlo, un
\emph{Essai} con la stessa idea; lo si conobbe solo nel 1813, quando Jacques Français ne
espose il contenuto negli \emph{Annales} di Gergonne invitando l'autore sconosciuto a
farsi avanti. A imporre la rappresentazione fu l'autorità di Gauss, che nel 1831 la difese
pubblicamente e coniò il nome di «numero complesso»: se $+1$, $-1$ e $\sqrt{-1}$, scrisse,
si fossero chiamate unità «diretta, inversa e laterale» anziché positiva, negativa e
immaginaria, nessuno vi avrebbe visto alcun mistero. È la lezione della
\cref{fig:prodotto-rotazione}: $i$ non è un numero misterioso, è una rotazione di un
angolo retto.
\end{notastorica}

\begin{esempio}[Dagli assi vento agli assi corda di un profilo]\label{es:assi-profilo}
Il significato «rotatorio» del prodotto è uno strumento di calcolo molto pratico. Un
profilo alare investito da una corrente di velocità $V_{\infty}$ con incidenza $\alpha$
è soggetto a una forza aerodinamica che nel riferimento \emph{vento} (asse $x$ nella
direzione della corrente, asse $y$ perpendicolare, verso l'alto) ha componenti $D$
(resistenza) e $L$ (portanza). Rappresentiamo la forza con il numero complesso
$F_{v}=D+iL$. Nel riferimento \emph{corda} (asse lungo la corda, dal bordo d'attacco al
bordo d'uscita, e asse normale), che con il profilo cabrato è ruotato di un angolo
$-\alpha$ rispetto al riferimento vento, le componenti si ottengono ruotando il vettore
in senso opposto, cioè moltiplicando per $\cos\alpha+i\sen\alpha$:
\[
  F_{c}=F_{v}(\cos\alpha+i\sen\alpha)=(D\cos\alpha-L\sen\alpha)+i\,(L\cos\alpha+D\sen\alpha).
\]
La parte reale è la \emph{forza assiale} $A=D\cos\alpha-L\sen\alpha$, il coefficiente
dell'immaginario è la \emph{forza normale} $N=L\cos\alpha+D\sen\alpha$: sono le formule
che si trovano in ogni testo di aerodinamica, ottenute qui con una sola moltiplicazione.
Con $L=1000$~N, $D=50$~N e $\alpha=5^{\circ}$ si trova $N\approx1000{,}6$~N e
$A\approx-37{,}3$~N: la forza assiale è
\emph{negativa}, cioè diretta verso il bordo d'attacco (è l'effetto della
«suzione» sul bordo d'attacco). Il passaggio inverso si ottiene moltiplicando per
$\cos\alpha-i\sen\alpha$, cioè dividendo: $F_{v}=F_{c}/(\cos\alpha+i\sen\alpha)$.
\end{esempio}

\begin{figure}[htb]
\centering
\begin{tikzpicture}[scale=1.0, every node/.style={font=\footnotesize},
  etich/.style={fill=white, inner sep=1pt, fill opacity=0.85, text opacity=1}]
  \def\alf{15}\def\Lv{2.6}\def\Dv{0.3}
  \pgfmathsetmacro{\Nv}{\Lv*cos(\alf)+\Dv*sin(\alf)}
  \pgfmathsetmacro{\Av}{\Dv*cos(\alf)-\Lv*sin(\alf)}
  % corrente
  \foreach \yy in {-0.7,-0.1,0.5} {\draw[-Latex, grigio!70] (-3.6,\yy) -- (-2.5,\yy);}
  \node[grigio] at (-3.05,1.0) {$V_{\infty}$};
  % profilo NACA 0012 cabrato di alfa, con il quarto di corda nell'origine
  \begin{scope}[rotate=-\alf, scale=6.5, shift={(-0.25,0)}]
    \filldraw[fill=bluchiaro, draw=blu, thick]
      plot[domain=0:1, samples=70, smooth] (\x, {0.6*(0.2969*sqrt(\x)-0.1260*\x-0.3516*\x^2+0.2843*\x^3-0.1015*\x^4)})
      -- plot[domain=1:0, samples=70, smooth] (\x, {-0.6*(0.2969*sqrt(\x)-0.1260*\x-0.3516*\x^2+0.2843*\x^3-0.1015*\x^4)}) -- cycle;
    \draw[grigio, dashed, thin] (-0.08,0) -- (1.25,0);
  \end{scope}
  % riferimento orizzontale e angolo di incidenza
  \draw[grigio, dotted] (0,0) -- (5.6,0);
  \draw[grigio] (4.6,0) arc (0:-\alf:4.6);
  \node[grigio] at ({4.85*cos(\alf/2)},{-4.85*sin(\alf/2)}) {$\alpha$};
  % risultante
  \draw[-Latex, grigio, thick] (0,0) -- (\Dv,\Lv) node[above=1pt, etich] {$R$};
  % assi vento: D, L
  \draw[-Latex, rame, very thick] (0,0) -- (0,\Lv) node[left=5pt, etich] {$L$};
  \draw[-Latex, rame, very thick] (0,0) -- (\Dv,0);
  \node[rame, etich] at (0.45,-0.28) {$D$};
  % assi corda: N, A
  \draw[-Latex, verde, very thick] (0,0) -- ({\Nv*sin(\alf)},{\Nv*cos(\alf)}) node[right=2pt, etich] {$N$};
  \draw[-Latex, verde, very thick] (0,0) -- ({\Av*cos(\alf)},{-\Av*sin(\alf)});
  \node[verde, etich] at (-1.05,0.75) {$A<0$};
  \draw[verde, thin] (-0.8,0.62) -- ({0.6*\Av*cos(\alf)},{-0.6*\Av*sin(\alf)+0.04});
  \fill (0,0) circle (1.3pt);
\end{tikzpicture}
\caption{Un profilo NACA 0012 a incidenza $\alpha$ (qui esagerata a $15^{\circ}$ per
leggibilità; forze non in scala). La risultante $R$ è disegnata nel quarto di corda, che
per un profilo simmetrico è, in prima approssimazione, il centro di pressione. La forza
aerodinamica ha componenti $(D,L)$ negli assi vento e $(A,N)$ negli assi corda: si
passa dalle une alle altre moltiplicando per
$\cos\alpha\pm i\sen\alpha$.}
\label{fig:profilo-assi}
\end{figure}

\subsection{Potenze e radici}

\begin{corollario}[Formula di De Moivre]\label{cor:de-moivre}\index{formula di De Moivre}
Se $z=\rho(\cos\vartheta+i\sen\vartheta)$, per ogni $n\in\N$
\begin{equation}\label{eq:de-moivre}
  z^{n}=\rho^{n}(\cos n\vartheta+i\sen n\vartheta).
\end{equation}
La formula vale anche per $n=0$ e, se $z\neq0$, per $n$ intero negativo.
\end{corollario}

\begin{schema}
\Ipotesi{$z=\rho(\cos\vartheta+i\sen\vartheta)$; \ $n\in\N$, oppure $n=0$, oppure
$n\in\Z$ negativo e $z\neq0$.}
\Tesi{$z^{n}=\rho^{n}(\cos n\vartheta+i\sen n\vartheta)$.}
\end{schema}

\begin{strategia}
Per $n\in\N$ si procede per induzione: per il \cref{teo:prodotto-trigonometrico}, ogni
moltiplicazione per $z$ moltiplica il modulo per $\rho$ e aggiunge $\vartheta$
all'argomento. Il caso $n=0$ si riduce alla convenzione $z^{0}=1$; per gli esponenti
negativi si scrive $z^{-m}=1/z^{m}$ e si usa la formula del quoziente.
\end{strategia}

\begin{dimostrazione}
Per $n=1$ è ovvia; se vale per $n$, per la \eqref{eq:prodotto-trigonometrico}
$z^{n+1}=z^{n}\cdot z=\rho^{n}\rho\bigl[\cos(n\vartheta+\vartheta)+i\sen(n\vartheta+\vartheta)\bigr]$,
che è la \eqref{eq:de-moivre} con $n+1$. Per $n=0$ entrambi i membri valgono $1$; per
$n=-m$ con $m\in\N$, si applica la \eqref{eq:quoziente-trigonometrico} a $1/z^{m}$.
\end{dimostrazione}

Un numero complesso $w$ si dice \emph{radice $n$-esima} di $z$ se risulta $w^{n}=z$.
Nel campo reale un numero positivo ha una sola radice $n$-esima positiva, e i numeri
negativi non hanno radici di ordine pari; nel campo complesso la situazione è molto più
regolare.

\begin{teorema}[Radici $n$-esime di un numero complesso]\label{teo:radici-complesse}\index{radice $n$-esima!di un numero complesso}
Sia $z=\rho(\cos\vartheta+i\sen\vartheta)\neq0$ e sia $n\in\N$. Allora $z$ ha esattamente
$n$ radici $n$-esime distinte, date da
\begin{equation}\label{eq:radici-complesse}
  w_{k}=\sqrt[n]{\rho}\,\Bigl(\cos\frac{\vartheta+2k\pi}{n}+i\sen\frac{\vartheta+2k\pi}{n}\Bigr),
  \qquad k=0,1,\dots,n-1.
\end{equation}
\end{teorema}

\begin{schema}
\Ipotesi{$z\in\C$, $z\neq0$, in forma trigonometrica con modulo $\rho>0$ e argomento
$\vartheta$; \ $n\in\N$.}
\Tesi{Le soluzioni dell'equazione $w^{n}=z$ sono esattamente gli $n$ numeri distinti
$w_{0},\dots,w_{n-1}$ della \eqref{eq:radici-complesse}.}
\end{schema}

\begin{strategia}
Cerchiamo $w$ in forma trigonometrica, $w=r(\cos\varphi+i\sen\varphi)$. Per la formula
di De Moivre l'equazione $w^{n}=z$ diventa un'uguaglianza tra due numeri in forma
trigonometrica, e due numeri complessi non nulli sono uguali se e solo se hanno lo
stesso modulo e argomenti che differiscono per un multiplo di $2\pi$. Si ottengono così
due condizioni separate: una sul modulo, $r^{n}=\rho$, che ha l'unica soluzione
$r=\sqrt[n]\rho$ per il \cref{teo:radice-n}; e una sull'argomento,
$n\varphi=\vartheta+2k\pi$, che ha infinite soluzioni $\varphi=(\vartheta+2k\pi)/n$, una
per ogni $k\in\Z$. Resta da contare quante di queste danno numeri \emph{diversi}: due
valori di $k$ danno lo stesso $w$ se e solo se gli argomenti differiscono per un multiplo
di $2\pi$, cioè se e solo se i due $k$ differiscono per un multiplo di $n$. I valori
$k=0,\dots,n-1$ sono quindi tutti e soli i rappresentanti distinti.
\end{strategia}

\begin{dimostrazione}
Poiché $z\neq0$, anche ogni radice $w$ è diversa da $0$ (se fosse $w=0$ sarebbe
$z=w^{n}=0$); possiamo quindi scrivere $w=r(\cos\varphi+i\sen\varphi)$, con $r>0$. Per la
\eqref{eq:de-moivre},
\begin{equation}\label{eq:radici-passo}
  w^{n}=z \quad\sse\quad r^{n}(\cos n\varphi+i\sen n\varphi)=\rho(\cos\vartheta+i\sen\vartheta).
\end{equation}
Prendendo i moduli, $r^{n}=\rho$, quindi $r=\sqrt[n]\rho$ (\cref{teo:radice-n}); allora
la \eqref{eq:radici-passo} equivale a $\cos n\varphi=\cos\vartheta$ e
$\sen n\varphi=\sen\vartheta$, cioè al fatto che gli angoli $n\varphi$ e $\vartheta$
individuino lo stesso punto della circonferenza di raggio $1$; ciò accade se e solo se
$n\varphi=\vartheta+2k\pi$ per qualche $k\in\Z$. Le radici sono dunque tutti e soli i
numeri della \eqref{eq:radici-complesse} con $k\in\Z$.

Siano $h,k\in\Z$. Gli argomenti $\frac{\vartheta+2h\pi}{n}$ e $\frac{\vartheta+2k\pi}{n}$
differiscono di $\frac{2(h-k)\pi}{n}$, che è un multiplo intero di $2\pi$ se e solo se
$h-k$ è un multiplo di $n$: quindi $w_{h}=w_{k}$ se e solo se $n$ divide $h-k$. In
particolare $w_{0},\dots,w_{n-1}$ sono distinti (due indici tra $0$ e $n-1$ differiscono
per meno di $n$, e non sono uguali). D'altra parte, ogni $k\in\Z$ si scrive come
$k=qn+s$ con $q=\lfloor k/n\rfloor\in\Z$ e $s\in\{0,\dots,n-1\}$
(\cref{prop:parte-intera}), e allora $w_{k}=w_{s}$. Le radici distinte sono quindi
esattamente $n$.
\end{dimostrazione}

In particolare il modulo di una radice $n$-esima è uguale alla radice $n$-esima
(aritmetica) del modulo, e gli argomenti delle $n$ radici sono equidistanziati di
$2\pi/n$: \emph{le radici $n$-esime di $z$ sono i vertici di un poligono regolare di $n$
lati inscritto nella circonferenza di centro $0$ e raggio $\sqrt[n]{\abs z}$}
(\cref{fig:radici-unita}).

\begin{esempio}[Le radici quadrate di $i$]\label{es:radici-i}
Calcoliamo le radici quadrate del numero complesso $z=i$. In forma trigonometrica $i$ ha
modulo uguale a $1$ e argomento uguale a $\pi/2$. Quindi, in base alla
\eqref{eq:radici-complesse}, le due radici quadrate di $i$ hanno modulo uguale a $1$ e
argomento uguale a
\begin{equation}\label{eq:argomenti-radici-i}
  \varphi_{0}=\frac{\pi/2}{2}=\frac\pi4, \qquad \varphi_{1}=\frac{\pi/2+2\pi}{2}=\frac54\pi .
\end{equation}
In corrispondenza si ottiene
\begin{align*}
  w_{0}&=\cos\frac\pi4+i\sen\frac\pi4=\frac{\sqrt2}{2}+i\,\frac{\sqrt2}{2},\\
  w_{1}&=\cos\frac54\pi+i\sen\frac54\pi=-\frac{\sqrt2}{2}-i\,\frac{\sqrt2}{2}.
\end{align*}
Quindi le due radici quadrate di $i$ sono
\begin{equation}\label{eq:radici-i}
  \pm\Bigl(\frac{\sqrt2}{2}+i\,\frac{\sqrt2}{2}\Bigr)=\pm\frac{\sqrt2}{2}(1+i).
\end{equation}
Verifica: $\bigl(\frac{\sqrt2}{2}(1+i)\bigr)^{2}=\frac12(1+2i+i^{2})=\frac12\cdot2i=i$.
\end{esempio}

\begin{figure}[htb]
\centering
\begin{tikzpicture}[scale=1.3, every node/.style={font=\footnotesize}]
  \draw[-Latex, grigio] (-1.45,0) -- (1.55,0) node[above left] {$\operatorname{Re}$};
  \draw[-Latex, grigio] (0,-1.35) -- (0,1.5) node[below left] {$\operatorname{Im}$};
  \draw[grigio!60] (0,0) circle (1);
  \draw[blu, thick, fill=bluchiaro, fill opacity=0.5] (0:1) \foreach \k in {1,...,5} { -- (\k*60:1)} -- cycle;
  \foreach \k in {0,...,5} {\fill[blu] (\k*60:1) circle (1.4pt);}
  \node[below right] at (0:1) {$w_{0}$};
  \node[above right] at (60:1) {$w_{1}$};
  \node[above left] at (120:1) {$w_{2}$};
  \node[below left] at (180:1) {$w_{3}$};
  \node[below left] at (240:1) {$w_{4}$};
  \node[below right] at (300:1) {$w_{5}$};
  \draw[rame] (0.35,0) arc (0:60:0.35); \node[rame] at (30:0.55) {$\frac{\pi}{3}$};
\end{tikzpicture}\hfill
\begin{tikzpicture}[scale=1.3, every node/.style={font=\footnotesize}]
  \draw[-Latex, grigio] (-1.45,0) -- (1.55,0) node[above left] {$\operatorname{Re}$};
  \draw[-Latex, grigio] (0,-1.35) -- (0,1.5) node[below left] {$\operatorname{Im}$};
  \draw[grigio!60] (0,0) circle (1);
  \fill[verde] (90:1) circle (1.5pt) node[above right] {$i$};
  \draw[rame, thick] (45:1) -- (225:1);
  \fill[rame] (45:1) circle (1.5pt) node[above right] {$w_{0}$};
  \fill[rame] (225:1) circle (1.5pt) node[below left] {$w_{1}$};
  \draw[verde, -Latex] (0.55,0) arc (0:90:0.55); \node[verde] at (68:0.75) {$\frac\pi2$};
  \draw[rame, -Latex] (0.3,0) arc (0:45:0.3); \node[rame] at (22:0.5) {$\frac\pi4$};
\end{tikzpicture}
\caption{A sinistra: le sei radici seste dell'unità $w_{k}=\cos\frac{k\pi}{3}+i\sen\frac{k\pi}{3}$
($w_{0}=1$, $w_{1}=\frac12+i\frac{\sqrt3}{2}$, $w_{3}=-1$, \dots), vertici di un esagono
regolare inscritto nella circonferenza unitaria. A destra: le due radici quadrate di $i$
(\cref{es:radici-i}), $w_{0}=\frac{\sqrt2}{2}(1+i)$ e $w_{1}=-w_{0}$, con argomenti $\pi/4$
e $5\pi/4$: il «poligono regolare di due lati» è un diametro.}
\label{fig:radici-unita}
\end{figure}

\subsection{Equazioni di secondo grado}

Possiamo ora mantenere la promessa iniziale: in $\C$ ogni equazione di secondo grado ha
soluzione.

\needspace{10\baselineskip}
\begin{proposizione}[Equazioni di secondo grado in $\C$]\label{prop:secondo-grado-C}
Siano $a,b,c\in\C$ con $a\neq0$, e sia $\Delta=b^{2}-4ac$. Detta $\delta$ una qualsiasi
radice quadrata complessa di $\Delta$ (che esiste per il \cref{teo:radici-complesse},
oppure è $\delta=0$ se $\Delta=0$), le soluzioni dell'equazione $az^{2}+bz+c=0$ sono
\begin{equation}\label{eq:soluzioni-C}
  z_{1,2}=\frac{-b\pm\delta}{2a},
\end{equation}
e valgono le relazioni \eqref{eq:viete}. In particolare, se $a,b,c$ sono reali e
$\Delta<0$, le soluzioni sono complesse coniugate:
$z_{1,2}=-\dfrac{b}{2a}\pm i\,\dfrac{\sqrt{-\Delta}}{2a}$.
\end{proposizione}

\begin{schema}
\Ipotesi{$a,b,c\in\C$, \ $a\neq0$; \ $\Delta=b^{2}-4ac$; \ $\delta\in\C$ con
$\delta^{2}=\Delta$.}
\Tesi{$az^{2}+bz+c=0\sse z=\frac{-b\pm\delta}{2a}$, \ e valgono le \eqref{eq:viete}; \
se $a,b,c\in\R$ e $\Delta<0$, \ $z_{1,2}=-\frac{b}{2a}\pm i\frac{\sqrt{-\Delta}}{2a}$.}
\end{schema}

\begin{strategia}
Si ripete il completamento del quadrato, che usa soltanto le proprietà di campo e
quindi vale in $\C$ senza modifiche. L'unico ingrediente nuovo è l'esistenza di una
radice quadrata $\delta$ di $\Delta$, garantita dal \cref{teo:radici-complesse} anche
quando $\Delta$ non è un reale non negativo. Scritto il trinomio come prodotto di due
fattori di primo grado, la legge di annullamento del prodotto conclude.
\end{strategia}

\begin{dimostrazione}
Il «completamento del quadrato» usa solo gli assiomi di campo, e vale quindi anche in
$\C$:
\[
  az^{2}+bz+c=a\Bigl[\Bigl(z+\frac{b}{2a}\Bigr)^{2}-\frac{b^{2}-4ac}{4a^{2}}\Bigr]
  =a\Bigl(z+\frac{b}{2a}-\frac{\delta}{2a}\Bigr)\Bigl(z+\frac{b}{2a}+\frac{\delta}{2a}\Bigr),
\]
dove abbiamo usato $\delta^{2}=\Delta$ e la scomposizione $u^{2}-v^{2}=(u-v)(u+v)$. Per
la legge di annullamento del prodotto (valida in $\C$, che è un campo) il primo membro è
nullo se e solo se uno dei due fattori lo è, cioè se e solo se $z$ è uno dei numeri
\eqref{eq:soluzioni-C}. Le relazioni \eqref{eq:viete} seguono sommando e moltiplicando:
$z_{1}+z_{2}=-b/a$ e $z_{1}z_{2}=(b^{2}-\delta^{2})/(4a^{2})=c/a$. Se infine $a,b,c$ sono
reali e $\Delta<0$, una radice quadrata di $\Delta$ è $\delta=i\sqrt{-\Delta}$, perché
$(i\sqrt{-\Delta})^{2}=i^{2}(-\Delta)=\Delta$.
\end{dimostrazione}

Ad esempio, l'equazione $z^{2}+1=0$ da cui eravamo partiti ha $\Delta=-4$, $\delta=2i$ e
soluzioni $z_{1,2}=\pm i$; e ora le relazioni $z_{1}+z_{2}=0$, $z_{1}z_{2}=1$ hanno
pienamente senso: $i+(-i)=0$ e $i\cdot(-i)=-i^{2}=1$.

\begin{notastorica}[Il teorema fondamentale dell'algebra]
La \cref{prop:secondo-grado-C} è il caso più semplice di un risultato molto più
generale, il \emph{teorema fondamentale dell'algebra}: \emph{ogni polinomio di grado
$n\ge1$ a coefficienti complessi ha almeno una radice in $\C$} (e quindi, scomponendo,
esattamente $n$ radici contate con la loro molteplicità). Enunciato già da Albert Girard
(1629) e tentato da d'Alembert (1746) e da Eulero, fu dimostrato in modo sostanzialmente
completo da Carl Friedrich Gauss nella sua tesi di dottorato (1799; la dimostrazione
conteneva una lacuna di natura topologica, colmata solo nel Novecento), e poi da
Jean-Robert Argand (1806, in una versione ripresa nel 1813), il dilettante incontrato
nella nota sul piano di Gauss, che fu anche il primo a enunciarlo esplicitamente per i
polinomi a coefficienti complessi. Il teorema dice che con $\C$ il processo di ampliamento iniziato passando da $\N$
a $\Z$ (per sottrarre), proseguito con $\Q$ (per dividere) e con $\R$ (per estrarre radici
e passare al limite) si \emph{chiude}: non occorrono altri numeri per risolvere le equazioni
algebriche.
\end{notastorica}

\subsection{La forma esponenziale}

Segnaliamo infine una notazione \emph{di tipo esponenziale} per i numeri complessi.
Posto
\begin{equation}\label{eq:eulero}
  e^{i\vartheta}=\cos\vartheta+i\sen\vartheta \qquad (\vartheta\in\R),
\end{equation}
un numero complesso $z=\rho(\cos\vartheta+i\sen\vartheta)$ si potrà rappresentare nella
forma
\begin{equation}\label{eq:forma-esponenziale}
  z=\rho\,e^{i\vartheta}.
\end{equation}
Per ora la \eqref{eq:eulero} è soltanto una \emph{definizione}, cioè un'abbreviazione:
nel capitolo sulle serie, dopo aver definito $e^{z}$ per ogni $z\in\C$ come somma della
serie esponenziale, vedremo che essa è in realtà un teorema, la \emph{formula di Eulero}
(\cref{teo:formula-eulero}), che lega la funzione esponenziale alle funzioni
trigonometriche. La notazione,
in ogni caso, è giustificata dal fatto che si comporta come una potenza. In tal modo,
infatti, la \eqref{eq:prodotto-trigonometrico} si riscrive come
\begin{equation}\label{eq:prodotto-esponenziale}
  z\cdot z'=\rho\,e^{i\vartheta}\cdot\rho'e^{i\vartheta'}=\rho\rho'\,e^{i(\vartheta+\vartheta')},
\end{equation}
che è la stessa regola $e^{u}e^{v}=e^{u+v}$ delle potenze reali; la formula di De Moivre
diventa $(\rho e^{i\vartheta})^{n}=\rho^{n}e^{in\vartheta}$, e le radici
\eqref{eq:radici-complesse} si scrivono
$w_{k}=\sqrt[n]\rho\,e^{i(\vartheta+2k\pi)/n}$. Se $w=r\,e^{i\varphi}$ è una radice
$n$-esima di $z=\rho\,e^{i\vartheta}$, sarà $r^{n}e^{in\varphi}=\rho\,e^{i\vartheta}$,
perciò $r=\sqrt[n]\rho$, $\cos n\varphi=\cos\vartheta$ e $\sen n\varphi=\sen\vartheta$:
ne seguono di nuovo le \eqref{eq:radici-complesse}.

\begin{notastorica}[La «formula più bella»]
Con $\vartheta=\pi$ la \eqref{eq:eulero} dà $e^{i\pi}=\cos\pi+i\sen\pi=-1$, cioè
\[
  e^{i\pi}+1=0,
\]
un'identità che lega le cinque costanti fondamentali della matematica: $0$, $1$, $e$,
$\pi$, $i$. Compare, nella forma $\cos x+i\sen x=e^{ix}$, nell'\emph{Introductio in
analysin infinitorum} di Leonhard Euler (1748); il simbolo $i$ per $\sqrt{-1}$ fu
introdotto dallo stesso Eulero nel 1777 e reso popolare da Gauss. In un sondaggio tra
matematici promosso dalla rivista \emph{The Mathematical Intelligencer} (1988) l'identità
di Eulero fu votata come «il più bel teorema della matematica».
\end{notastorica}

\begin{idea}[Oscillazioni e fasori]
Per un ingegnere i numeri complessi sono soprattutto il linguaggio delle
\emph{oscillazioni}. Uno spostamento armonico, come la flessione all'estremità di un'ala
che vibra, $x(t)=X\cos(\omega t+\phi)$, si scrive
$x(t)=\operatorname{Re}\bigl(X e^{i\phi}\,e^{i\omega t}\bigr)$: il numero complesso
$X e^{i\phi}$ (il \emph{fasore}) contiene ampiezza e fase, e il fattore $e^{i\omega t}$
lo fa ruotare nel piano di Gauss con velocità angolare $\omega$. Sommare due oscillazioni
della stessa frequenza diventa una somma di vettori (\cref{fig:piano-gauss}), derivare
rispetto al tempo diventa moltiplicare per $i\omega$, cioè ruotare di $90^{\circ}$ e
dilatare. È così che si studiano le vibrazioni strutturali, il flutter aeroelastico e i
circuiti in corrente alternata.
\end{idea}

\begin{python}[Numeri complessi in Python]
import cmath, math

z = 3 + 2j          # in Python l'unità immaginaria si scrive j
w = 1 - 1j
print(z / w)                       # quoziente: (0.5+2.5j)
print(abs(z), z.conjugate())       # modulo e coniugato
print(cmath.polar(1j))             # (modulo, argomento) di i

# Radici n-esime con la formula del teorema
def radici(z, n):
    rho, theta = cmath.polar(z)
    r = rho ** (1 / n)
    return [cmath.rect(r, (theta + 2 * k * math.pi) / n)
            for k in range(n)]

for w_k in radici(1j, 2):
    print(f"{w_k:.6f}   (w_k)^2 = {w_k**2:.6f}")

# Dagli assi vento agli assi corda (esempio): F_c = F_v * e^(i alfa)
L, D, alfa = 1000.0, 50.0, math.radians(5)
Fc = complex(D, L) * cmath.exp(1j * alfa)
print(f"A = {Fc.real:.1f} N,  N = {Fc.imag:.1f} N")

# Il calcolo di Bombelli: (2 + i)^3 = 2 + 11i, cioè 2 + sqrt(-121)
print((2 + 1j) ** 3)
\end{python}
\begin{risultato}
(0.5+2.5j)
3.605551275463989 (3-2j)
(1.0, 1.5707963267948966)
0.707107+0.707107j   (w_k)^2 = 0.000000+1.000000j
-0.707107-0.707107j   (w_k)^2 = 0.000000+1.000000j
A = -37.3 N,  N = 1000.6 N
(2+11j)
\end{risultato}

\begin{figure}[htb]
\centering
\includegraphics[width=\textwidth]{radici_unita}
\caption{Le radici $n$-esime dell'unità per $n=3,\dots,8$, calcolate con la
\eqref{eq:radici-complesse} (\texttt{python/fig\_radici\_unita.py}): al crescere di $n$
i poligoni regolari approssimano sempre meglio la circonferenza.}
\label{fig:radici-unita-python}
\end{figure}
